\documentclass[11pt,a4paper]{article}

\usepackage[utf8]{inputenc}
\usepackage[T1]{fontenc}
\usepackage{times}
\usepackage{helvet}
\usepackage{courier}
\usepackage{fullpage}
\usepackage{amsmath,amssymb,amsfonts}
\usepackage{graphicx}
\usepackage{booktabs}
\usepackage{multirow}
\usepackage{hyperref}
\usepackage{listings}
\usepackage{xcolor}
\usepackage{algorithm}
\usepackage{algorithmic}
\usepackage{natbib}

\lstset{
    basicstyle=\ttfamily\scriptsize,
    numbers=left,
    numberstyle=\tiny,
    stepnumber=1,
    numbersep=5pt,
    backgroundcolor=\color{white},
    showspaces=false,
    showstringspaces=false,
    frame=single,
    breaklines=true,
    postbreak=\mbox{\hookrightarrow\space},
}

\hypersetup{
    colorlinks=true,
    linkcolor=blue,
    citecolor=blue,
    urlcolor=blue,
}

\title{\textbf{SSLM: Achieving Useful Agent Capabilities at GPT-2 Scale} \\
\Large{A Super Small But Usable Architecture}}

\author{
    Zice Wang \\
    \texttt{wangzc1@mails.neu.edu.cn} \\
    NTAIRS, Northeastern University, China \\
}

\date{May 21, 2026}

\begin{document}

\maketitle

\begin{abstract}
We present SSLM, a hybrid decoder-only language model designed to achieve \textbf{useful agent-level capabilities at GPT-2 scale (~1.5B parameters)}. The architecture combines three proven innovations: Multi-head Latent Attention (MLA) from DeepSeek V2/V3 for compressed KV-cache and multi-turn conversation memory efficiency, Mamba2 state-space layers for efficient long-context processing and agent planning, and an aux-loss-free Mixture-of-Experts (MoE) module for parameter-efficient scaling without load-balancing auxiliary losses. SSLM v2 achieves \textbf{1.45B total parameters} (comparable to GPT-2 XL) with only $\sim$768M active parameters per forward pass, making it well-suited for efficient agent inference. We present complete training results from 22K vocab (SSLM v1) to 53K vocab (SSLM v2), alongside evaluations on C-Eval, TruthfulQA, and GSM8K benchmarks.
\end{abstract}

\section{Introduction}
\label{sec:introduction}

\subsection{Goal: Useful Agent at GPT-2 Scale}

The original GPT-2 models \citep{gpt2} (117M / 345M / 774M / 1.5B) established that decoder-only transformers at the $\sim$1B parameter scale can exhibit surprising language understanding and generation capabilities. Our goal is to match or exceed GPT-2's utility while being more parameter-efficient, by combining modern architectural innovations (MLA, Mamba2, MoE) that were not available in 2019.

\textbf{``Useful agent level''} means the model should be capable of:

\begin{enumerate}
    \item \textbf{Tool use:} Calling external tools/APIs in a coherent sequence (ReAct-style) \citep{react}
    \item \textbf{Multi-step reasoning:} Breaking down complex tasks into subgoals and executing them
    \item \textbf{Long-horizon memory:} Maintaining context across extended conversations ($\ge$ 1024 tokens)
    \item \textbf{Multilingual tool navigation:} Operating on both English and Chinese inputs
    \item \textbf{Code generation \& execution:} Producing correct code for simple tasks
\end{enumerate}

This goal differentiates DenseSLM4MoE from pure benchmark-optimized models --- we prioritize architectural choices that enable agent workflows, not just high C-Eval scores.

\subsection{Why Hybrid MLA + Mamba2 + MoE?}

For agent workloads, three capabilities matter:

\begin{table}[h]
    \centering
    \caption{Architectural solutions for agent capabilities}
    \begin{tabular}{lll}
        \toprule
        \textbf{Capability} & \textbf{Requirement} & \textbf{Solution} \\
        \midrule
        Long context memory & KV-cache grows with turns & MLA: compresses KV to 512-dim latent \\
        Efficient processing & 48 layers needed & Mamba2: O(N) SSM for 75\% of layers \\
        Parameter scale & 1B+ params, manageable compute & MoE: 64 experts, top-2 active \\
        \bottomrule
    \end{tabular}
\end{table}

\subsection{Contributions}

\begin{enumerate}
    \item \textbf{Agent-oriented hybrid architecture:} 3:1 Mamba2:MLA layer pattern --- Mamba2 layers handle efficient sequence processing; MLA layers provide high-quality attention with compressed KV-cache for multi-turn memory.
    \item \textbf{Aux-loss-free MoE routing:} No load-balancing auxiliary loss; instead, a learnable \texttt{e\_score\_correction\_bias} updated via simple load-based feedback, inspired by DeepSeek V3 \citep{deepseekv3} and Nemotron-H \citep{nemotron}.
    \item \textbf{Grouped top-k routing:} $n\_group=8$, $topk\_group=2$ for coarse-to-fine expert selection, reducing routing decision overhead.
    \item \textbf{Complete agent-focused training pipeline:} Muon optimizer \citep{muon}, WCL learning rate schedule, multilingual pretraining on English/Chinese/code data, SwanLab tracking.
    \item \textbf{Ablation study} showing that wider experts (8 experts, top-2) outperform narrower experts (32 experts, top-8) at equal parameter count for our target tasks.
\end{enumerate}

\section{Architecture}
\label{sec:architecture}

\subsection{Design Rationale for Agent Use Cases}

For agent workloads, we made the following architectural decisions:

\textbf{Why 3:1 Mamba2:MLA?}
\begin{itemize}
    \item Mamba2 handles the bulk of forward pass computation (75\% of layers) at O(N) complexity, making long-context generation efficient.
    \item MLA layers appear every 4th layer to provide full attention capability where needed for agentic decisions (e.g., tool selection, instruction following).
    \item MoE only replaces FFN in MLA layers, keeping active compute low.
\end{itemize}

\textbf{Why compressed KV-cache via MLA?}
\begin{itemize}
    \item Agent systems need to maintain context across many turns. With standard MHA, the KV-cache grows as $2 \times num\_layers \times num\_kv\_heads \times seq\_len \times head\_dim$.
    \item MLA compresses all KV to a low-rank latent (512-dim), reducing cache memory by $\sim$10$\times$ for the same effective context.
    \item This enables serving longer agent conversations within the same memory budget.
\end{itemize}

\textbf{Why aux-loss-free MoE?}
\begin{itemize}
    \item Standard MoE with auxiliary load-balancing loss introduces gradient interference between the main language modeling objective and the routing objective.
    \item Aux-loss-free routing avoids this interference, leading to better loss convergence.
    \item The bias correction mechanism naturally balances expert utilization without destabilizing training.
\end{itemize}

\subsection{Overall Structure}

\begin{figure}[h!]
    \centering
    \rule{0.8\textwidth}{0.3\textwidth}
    \caption{DenseSLM4MoE Architecture: Embedding $\rightarrow$ 36 Mamba2 layers + 12 MLA+MoE layers $\rightarrow$ Final LayerNorm $\rightarrow$ LM Head}
    \label{fig:architecture}
\end{figure}

The model follows: Embedding $\rightarrow$ [Layer 0-2: Mamba2 $\times$ 3] $\rightarrow$ [Layer 3: MLA + MoE] $\rightarrow$ [Layer 4-6: Mamba2 $\times$ 3] $\rightarrow$ ... (pattern repeats 12 times) $\rightarrow$ Final LayerNorm $\rightarrow$ LM Head (tied to embedding).

\subsection{Mamba2 Layers}

Mamba2 \citep{mamba2} (from \texttt{fla} / \texttt{mamba\_ssm}) uses a hardware-aware chunk-scan algorithm for linear-time sequence modeling:

\begin{algorithm}[H]
\caption{Mamba2 Layer Forward}
\begin{algorithmic}
\STATE Input $\rightarrow$ RMSNormGated $\rightarrow$ [Chunk-scan SSM with state\_size=64, expand=2]
\STATE $\rightarrow$ Selective State Update $\rightarrow$ Conv1D (kernel=4) $\rightarrow$ Output
\end{algorithmic}
\end{algorithm}

\begin{table}[ht!]
    \centering
    \caption{Mamba2 Key Parameters}
    \begin{tabular}{lll}
        \toprule
        \textbf{Parameter} & \textbf{Value} & \textbf{Agent Implication} \\
        \midrule
        mamba\_expand & 2 & 2$\times$ hidden size expansion for SSM state \\
        mamba\_state\_size & 64 & Small state --- fast inference, limited memory \\
        mamba\_conv\_kernel & 4 & Short-term local context before SSM \\
        mamba\_n\_groups & 4 & Grouped state branching for expressiveness \\
        \bottomrule
    \end{tabular}
\end{table}

\subsection{Multi-head Latent Attention (MLA)}

MLA from DeepSeek V3 \citep{deepseekv3} compresses the KV cache:

\begin{align}
    Q &: hidden \rightarrow q\_lora\_rank(32) + qk\_nope(128 \times 16) + qk\_rope(64 \times 16) \\
    KV &: hidden \rightarrow kv\_lora\_rank(512) \quad \text{[compressed representation]}
\end{align}

\textbf{Key advantage for agents:} The compressed KV representation (512 dim) vs standard KV ($4 \times 128 \times 2 = 1024$ dim) saves $\sim$50\% KV cache memory per layer.

\subsection{MoE Module}

Each MLA layer's FFN is replaced by:

\textbf{Router} --- Single linear: $hidden(768) \rightarrow n\_routed\_experts(64)$, sigmoid scores + correction bias

\textbf{Routed Experts (64):}
\begin{itemize}
    \item Each expert: $hidden(768) \rightarrow moe\_intermediate(1024) \rightarrow hidden(768) = \sim$1.57M params
    \item Top-2 per token $\rightarrow$ $\sim$3.14M active params per MLA layer
    \item Total MoE params: $64 \times 2 \times 768 \times 1024 \times 2 \approx 201$M (across all 12 MLA layers)
\end{itemize}

\textbf{Shared Expert (always active):}
\begin{itemize}
    \item $hidden(768) \rightarrow shared\_intermediate(1024) \rightarrow hidden(768) = \sim$1.57M params
    \item Always active, no routing decision needed
\end{itemize}

\textbf{Grouped routing:}
\begin{equation}
    Scores \rightarrow reshape(-1, 8, 8) \rightarrow topk(2) \text{ per group} \rightarrow sum \rightarrow topk(2) groups \rightarrow mask \rightarrow topk(2) experts
\end{equation}

\subsection{Aux-Loss-Free Routing Algorithm}

During training, after each forward pass (executed in \texttt{update\_expert\_bias()}):

\begin{algorithm}[H]
\caption{Aux-Loss-Free Routing Bias Update}
\begin{algorithmic}
\STATE $load \leftarrow$ bincount(topk\_indices.reshape(-1), minlength=n\_routed\_experts)
\STATE $target \leftarrow load.mean()$
\STATE $error \leftarrow sign(target - last\_expert\_load)$
\STATE $e\_score\_correction\_bias \leftarrow e\_score\_correction\_bias + update\_rate \times error$
\STATE $e\_score\_correction\_bias \leftarrow e\_score\_correction\_bias - e\_score\_correction\_bias.mean()$ \COMMENT{zero-sum constraint}
\end{algorithmic}
\end{algorithm}

\textbf{Key properties:}
\begin{itemize}
    \item No gradient through the routing balance signal
    \item Biases evolve smoothly during training
    \item Zero-sum constraint ensures biases don't shift the overall score distribution
    \item $update\_rate = 1e-3$ (modest, gradual correction)
\end{itemize}

\subsection{Active vs Total Parameters}

\begin{table}[ht!]
    \centering
    \caption{Parameter Efficiency: Active vs Total}
    \begin{tabular}{lll}
        \toprule
        \textbf{Metric} & \textbf{Value} & \textbf{Explanation} \\
        \midrule
        Total parameters & 1.45B & Full model including all 64 experts \\
        Active parameters (per forward pass) & $\sim$768M & 2 experts $\times$ 12 MLA layers + shared expert + all Mamba2/MLA layers \\
        Embedding/LM head & 40.6M & Tied, not doubled \\
        MoE experts (stored, not active) & $\sim$201M & All 64 expert pairs stored \\
        \bottomrule
    \end{tabular}
\end{table}

This gives an \textbf{active ratio of $\sim$53\%}, meaning the model uses $\sim$768M FLOPs per token vs a hypothetical 1.45B for a dense model of the same size.

\section{Parameter Counts}
\label{sec:parameters}

\subsection{DenseSLM4MoE v2 (1.45B total)}

\textbf{Configuration:} vocab=52,877, hidden=768, layers=48, heads=16, MoE=64 experts top-2

\begin{table}[h]
    \centering
    \caption{DenseSLM4MoE v2 Parameter Breakdown}
    \begin{tabular}{llr}
        \toprule
        \textbf{Component} & \textbf{Formula} & \textbf{Count} \\
        \midrule
        Embedding & 52,877 $\times$ 768 & 40.6M \\
        Mamba2 $\times$ 36 & $\sim$2.5M/layer $\times$ 36 & $\sim$90M \\
        MLA $\times$ 12 & $\sim$5M/layer $\times$ 12 & $\sim$60M \\
        Router (gate) & 768 $\times$ 64 & 49K \\
        MoE routed experts & 64 $\times$ 2 $\times$ 768 $\times$ 1024 & 201M \\
        MoE shared expert $\times$ 12 & 2 $\times$ 768 $\times$ 1024 & 19M \\
        Final LayerNorm & 2 $\times$ 768 & 1.5K \\
        \midrule
        \textbf{Total} & & \textbf{$\sim$1,418M} (reported: 1,450.86M) \\
        \bottomrule
    \end{tabular}
\end{table}

\subsection{DenseSLM4 (Dense, $\sim$105M)}

\begin{table}[h]
    \centering
    \caption{DenseSLM4 Parameter Breakdown}
    \begin{tabular}{lr}
        \toprule
        \textbf{Component} & \textbf{Count} \\
        \midrule
        Embedding & 4.73M \\
        Mamba2 $\times$ 24 & $\sim$60M \\
        MLA $\times$ 8 + DenseMLP & $\sim$40M \\
        Final LayerNorm & $\sim$0.001M \\
        \midrule
        \textbf{Total} & \textbf{$\sim$105M} \\
        \bottomrule
    \end{tabular}
\end{table}

\subsection{Ablation Study: Expert Width vs Count}

Three configurations at $\sim$25M parameters (small Transformer backbone):

\begin{table}[h]
    \centering
    \caption{Expert Width Ablation Study}
    \begin{tabular}{clllrrr}
        \toprule
        \textbf{Config} & \textbf{Expert Dim} & \textbf{\# Experts} & \textbf{Top-K} & \textbf{Eval Loss} & \textbf{PPL} & \textbf{Params} \\
        \midrule
        A (narrow) & 192 & 32 & 8 & 2.113 & 8.27 & 24.83M \\
        \textbf{B (wide)} & \textbf{768} & \textbf{8} & \textbf{2} & \textbf{2.033} & \textbf{7.63} & \textbf{24.80M} \\
        C (medium) & 384 & 16 & 4 & 2.053 & 7.80 & 24.81M \\
        \bottomrule
    \end{tabular}
\end{table}

\textbf{Conclusion:} Wider experts with lower top-k consistently outperform narrower experts with higher top-k at equal parameter count. This informed our choice of $n\_routed\_experts=64, topk=2$.

\section{Training Infrastructure}
\label{sec:training}

\subsection{Optimizer: Muon + AdamW Hybrid}

\textbf{Muon} \citep{muon} (MomentUm Orthogonalized by Newton-Schulz):
\begin{itemize}
    \item Applied to: All linear weight matrices (Mamba2, MLA, MoE experts)
    \item Uses 5th-order Newton-Schulz iteration to orthogonalize gradient updates
    \item \texttt{moonlight} scaling mode: $0.2 \times \sqrt{max(A, B)} + wd \times 0.1$
\end{itemize}

\textbf{AdamW} for: Embeddings, biases, LayerNorms, LM head

\textbf{Why Muon for agent models?}
\begin{itemize}
    \item Muon achieves faster convergence on transformer weight matrices compared to vanilla AdamW in preliminary experiments
    \item Faster convergence means fewer training steps to reach the same capability level
\end{itemize}

\subsection{Learning Rate Schedule: WCL}

\textbf{WCL (Warmup $\rightarrow$ Cosine $\rightarrow$ Linear):}

\begin{align}
    LR = base\_lr \times \begin{cases}
        step / warmup & \text{if } step < warmup \\
        0.5 \times (1 + \cos(\pi \times (step - warmup) / (0.5 \times total - warmup))) & \text{if } step < 0.5 \times total \\
        1 - (step - 0.5 \times total) / (0.5 \times total) & \text{if } step \ge 0.5 \times total
    \end{cases}
\end{align}

For DenseSLM4MoE v2: warmup=9,000, total=190,967 $\rightarrow$ cosine phase from step 9,000 to $\sim$104,500.

\subsection{Agent-Specific Training Considerations}

For agent workloads, we made the following training choices:

\begin{enumerate}
    \item \textbf{Strided chunking (stride=896, seq\_len=1024):} Creates overlapping 128-token windows between chunks, simulating multi-turn agent conversations where context bleeds across turns.
    \item \textbf{Multilingual data (English + Chinese + Code):} Agent systems must operate across languages and interact with code.
    \item \textbf{Long sequence length (1024):} Chosen to support agent context windows of 10+ turns $\times$ $\sim$100 tokens/turn.
\end{enumerate}

\subsection{Training Speed}

\begin{itemize}
    \item \textbf{Platform:} NEU-lab2 GPU server (ssh-coordinator skill)
    \item \textbf{DenseSLM4MoE v2 speed:} $\sim$1.3s/step at batch\_size=32, seq\_len=1024
    \item \textbf{TF32} CUDA matmul enabled
    \item \textbf{Estimated total training time:} $\sim$70 hours for full 190K steps at current speed
\end{itemize}

\section{Dataset}
\label{sec:dataset}

\subsection{Composition (Agent-Focused)}

\begin{table}[h]
    \centering
    \caption{Pretraining Dataset Composition}
    \begin{tabular}{llrl}
        \toprule
        \textbf{Dataset} & \textbf{Domain} & \textbf{Samples} & \textbf{Agent Relevance} \\
        \midrule
        finepdfs\_edu\_50BT...fineweb\_edu\_20BT-shuffled & English web/edu/pdf & 2,502,288 & General instruction following \\
        the-stack-smol & GitHub code & 300,000 & Tool use, code generation \\
        Fineweb-Edu-Chinese-V2.1 & Chinese edu web & 717,776 & Chinese language agent tasks \\
        tiny-textbooks & Textbook QA & 399,000 & Multi-step reasoning \\
        orca\_math\_qa & Math word problems & 200,035 & Chain-of-thought reasoning \\
        Extra (2.parquet, 3.parquet) & Mixed & $\sim$16,476 & Diverse edge cases \\
        \midrule
        \textbf{Total base} & & \textbf{$\sim$4,135,575} & \\
        \bottomrule
    \end{tabular}
\end{table}

\begin{itemize}
    \item After chunking (stride=896, seq\_len=1024): $\sim$6.11M training samples
    \item Train/Val split: 99.9\% / 0.1\%
\end{itemize}

\subsection{Tokenization}

\begin{table}[h]
    \centering
    \caption{Tokenizer Versions}
    \begin{tabular}{lrll}
        \toprule
        \textbf{Tokenizer} & \textbf{Vocab Size} & \textbf{Used In} & \textbf{Compression Ratio} \\
        \midrule
        v1 & 22,232 & DenseSLM4MoE v1 & $\sim$1.3 tokens/char \\
        v2 & 52,877 & DenseSLM4MoE v2 & $\sim$1.5 tokens/char \\
        \bottomrule
    \end{tabular}
\end{table}

\subsection{Hybrid BPE Tokenizer Merging: Qwen3-8B + Custom BPE}

The v2 tokenizer leverages a \textbf{clever hybrid merging strategy} that combines the best of both worlds: the well-optimized Qwen3-8B BPE tokenizer with our custom-trained BPE tokenizer.

\textbf{Problem:} Qwen3-8B tokenizer has 151,669 tokens and 100K+ merge rules, but we only need a subset for our specific data distribution. Retraining from scratch loses the well-optimized English/Chinese/Code subword units learned by Qwen.

\textbf{Solution:} \texttt{rebuild\_tokenizer\_workspace.py} implements an intelligent merging pipeline:

\begin{enumerate}
    \item \textbf{Sample} text from the same 6-dataset training mixture used for pretraining
    \item \textbf{Count Qwen tokens}: Tokenize 100K samples with Qwen3-8B tokenizer, find the 49K most frequent non-special tokens
    \item \textbf{Recursive merge tracing}: For each selected token, recursively trace which BPE merge rules are needed. If token ``hello'' needs merges \texttt{["he", "llo"]}, we keep both sub-tokens and their constituent merges.
    \item \textbf{Merge}: Combine the pruned Qwen tokenizer with our custom-trained BPE (which learned domain-specific tokens like Chinese words, code patterns)
    \item \textbf{Result}: 52,877 vocab, $\sim$1.5 tokens/char compression, with only the necessary merge rules
\end{enumerate}

\begin{figure}[h]
    \centering
    \rule{0.8\textwidth}{0.2\textwidth}
    \caption{Hybrid Tokenizer Merging: Qwen3-8B (151K tokens) $\rightarrow$ Top-49K by frequency + Custom BPE $\rightarrow$ Merged 52K tokenizer}
    \label{fig:tokenizer}
\end{figure}

\textbf{Why this works:}
\begin{itemize}
    \item Qwen3-8B already learned excellent English/Chinese subword units from massive pretraining
    \item We only keep merges for tokens that actually appear in our data (data-aware pruning)
    \item Custom BPE adds domain-specific tokens (math notation, code patterns) that Qwen might under-represent
    \item The merged tokenizer achieves better compression ($\sim$1.5 tok/char) than either tokenizer alone on our specific mix
\end{itemize}

\textbf{Note:} \texttt{1.parquet} (Traditional Chinese) was excluded due to poor compression ratio (0.82, meaning 1 char $\approx$ 1 token), which is inefficient for Chinese agent tasks.

\subsection{Tokenizer Benchmark: Full Comparison}

We evaluate our tokenizer against four baselines across different vocabulary scales:
\begin{itemize}
    \item \textbf{GPT-2} (vocab=50,257) -- same-scale baseline
    \item \textbf{DeepSeek-V2} (vocab=100,000) -- 2x scale reference
    \item \textbf{DeepSeek-V3} (vocab=128,000) -- 2.5x scale reference
    \item \textbf{Qwen3-8B} (vocab=151,643) -- 3x scale reference
    \item \textbf{Ours} (vocab=52,877) -- our hybrid merged tokenizer
\end{itemize}

\begin{table}[h]
    \centering
    \caption{Tokenizer Compression Benchmark: tok/char (lower is better)}
    \resizebox{\textwidth}{!}{
    \begin{tabular}{llrrrrr}
        \toprule
        \textbf{Test} & \textbf{Source} & \textbf{GPT-2} & \textbf{DS-V2} & \textbf{DS-V3} & \textbf{Qwen3} & \textbf{Ours} \\
        \midrule
        English prose & Pride \& Prejudice (Gutenberg) & 0.2630 & 0.2425 & 0.2358 & \textbf{0.2357} & 0.2400 \\
        Code & Linux Makefile (GitHub) & 0.3817 & 0.3218 & 0.3178 & \textbf{0.2815} & 0.3123 \\
        Numbers & 10K random integers & \textbf{0.4006} & 1.0000 & 0.4350 & 1.0000 & 1.0000 \\
        Random chars & 10K random a-z & 0.5969 & 0.5364 & \textbf{0.5355} & 0.5401 & 0.5898 \\
        Random bytes & 5KB binary & 0.7882 & 0.8205 & 0.6514 & \textbf{0.6250} & 0.7185 \\
        Chinese (Simplified) & Dream of Red Chamber (converted) & 2.0971 & 0.8516 & \textbf{0.7327} & 0.7878 & 0.9045 \\
        Chinese (Traditional) & Dream of Red Chamber (Gutenberg) & 2.1442 & 1.1191 & \textbf{0.7791} & 0.8596 & 1.0087 \\
        \bottomrule
    \end{tabular}}
\end{table}

\textbf{Test file URLs (publicly accessible):}
\begin{itemize}
    \item English: \url{https://www.gutenberg.org/files/1342/1342-0.txt}
    \item Code: \url{https://raw.githubusercontent.com/torvalds/linux/master/Makefile}
    \item Chinese (Simplified): converted via OpenCC from Gutenberg source
    \item Chinese (Traditional): \url{https://www.gutenberg.org/cache/epub/24264/pg24264.txt}
\end{itemize}

\textbf{Key findings:}
\begin{itemize}
    \item \textbf{English:} Qwen3 (0.236) and DS-V3 (0.236) lead; Ours (0.240) is competitive despite 3x smaller vocab
    \item \textbf{Code:} Ours (0.312) and Qwen3 (0.282) lead; DeepSeek-V3 (0.318) close behind
    \item \textbf{Chinese (Simplified):} DS-V3 wins (0.73); Ours (0.90) is 23\% worse but still 57\% better than GPT-2
    \item \textbf{Chinese (Traditional):} DS-V3 wins (0.78); Ours (1.01) is 29\% worse but 53\% better than GPT-2
    \item \textbf{Numbers:} GPT-2 wins (0.40); all large-vocab tokenizers split digits granularly. Acceptable under agent tool-use paradigm with structured CoT.
    \item \textbf{Random bytes:} DS-V3 (0.65) and Qwen3 (0.63) lead
\end{itemize}

\textbf{Comparison with 128K-scale tokenizers:}
\begin{itemize}
    \item Ours (53K) vs DeepSeek-V3 (128K): English and Code are competitive (0.24 vs 0.24, 0.31 vs 0.32); Chinese gap is 23-29\%
    \item Ours (53K) vs Qwen3 (151K): Code gap on Linux Makefile (0.31 vs 0.28); English on par; Chinese gap 15-17\%
\end{itemize}

\textbf{Remarkable result:} Despite using only 53K vocabulary (vs 128-151K for reference tokenizers), Ours achieves competitive performance on English and Code. This demonstrates that \textbf{merging strategy quality > vocabulary size} in most domains except specialized Chinese text.

\subsection{Vocabulary Efficiency: $\ln(V)$ Scaling}

Raw tok/char comparisons favor larger vocabularies by design. To compare tokenizers fairly across different vocabulary sizes, we introduce an efficiency metric grounded in information theory.

\textbf{Theoretical basis:} BPE is a greedy approximation of Huffman coding. Under optimal encoding, each token carries $\log_2(V)$ bits of information, so the expected token count scales as $\text{tok/char} \approx A / \ln(V)$ where $A$ is a constant that captures how efficiently the vocabulary budget is used.

\textbf{Efficiency metric:} $A = \text{tok/char} \times \ln(V)$ \quad (lower $A$ is better)

A smaller $A$ means the tokenizer achieves better compression for the same vocabulary size, or equivalently, achieves the same compression with a smaller vocabulary.

\begin{table}[h]
    \centering
    \caption{Vocabulary-Normalized Efficiency: $A = \text{tok/char} \times \ln(V)$ (English, Pride \& Prejudice)}
    \begin{tabular}{lrrrr}
        \toprule
        \textbf{Tokenizer} & \textbf{Vocab ($V$)} & \textbf{$\ln(V)$} & \textbf{tok/char} & \textbf{$A$} \\
        \midrule
        \textbf{Ours} & 52,877 & 10.88 & 0.2400 & \textbf{2.610} \\
        DeepSeek-V3 & 128,000 & 11.76 & 0.2358 & 2.773 \\
        DeepSeek-V2 & 100,000 & 11.51 & 0.2425 & 2.791 \\
        Qwen3-8B & 151,643 & 11.93 & 0.2357 & 2.812 \\
        GPT-2 & 50,257 & 10.82 & 0.2630 & 2.845 \\
        \bottomrule
    \end{tabular}
\end{table}

\textbf{Vocabulary efficiency across test types ($A = \text{tok/char} \times \ln(V)$):}

\begin{table}[h]
    \centering
    \resizebox{\textwidth}{!}{
    \begin{tabular}{lrrrrr}
        \toprule
        \textbf{Test} & \textbf{GPT-2} & \textbf{DS-V2} & \textbf{DS-V3} & \textbf{Qwen3} & \textbf{Ours} \\
        \midrule
        English & 2.845 & 2.791 & 2.773 & 2.812 & \textbf{2.610} \\
        Code & 4.130 & 3.704 & 3.737 & \textbf{3.358} & 3.399 \\
        Chinese (Simplified) & 22.690 & 9.802 & \textbf{8.617} & 9.427 & 9.841 \\
        Chinese (Traditional) & 23.200 & 12.881 & \textbf{9.162} & 10.255 & 10.975 \\
        \bottomrule
    \end{tabular}}
\end{table}

\textbf{Key findings:}
\begin{itemize}
    \item Under $\ln(V)$ normalization, Ours achieves the lowest $A$ on English and Code -- \textbf{Ours is the most vocabulary-efficient tokenizer}
    \item On Chinese, DS-V3 leads (8.6-9.2) due to dedicated Chinese character optimization; Ours (9.8-11.0) is second among 50K-scale tokenizers
    \item GPT-2 has the highest $A$ (worst efficiency) despite having the smallest vocabulary -- confirming that vocabulary budget alone does not determine quality
    \item The $A$ values cluster in a narrow range (2.6-2.9 for English), consistent with the theoretical prediction that $A$ should be approximately constant across well-trained BPE tokenizers
\end{itemize}

\textbf{Theoretical interpretation:} The near-constancy of $A$ across tokenizers of different scales ($V$ from 50K to 151K) validates the $\ln(V)$ scaling law. Deviations from the mean reflect real differences in tokenizer quality. Ours achieves the best $A$ by combining Qwen3-8B's well-learned subword merges with custom BPE training on domain-specific data.

\subsection{Code Tokenization Deep Dive}

We analyze tokenization quality on Python quicksort code:

\begin{table}[h]
    \centering
    \caption{Python Quicksort: Token Count Comparison}
    \begin{tabular}{lrrr}
        \toprule
        \textbf{Tokenizer} & \textbf{Vocab} & \textbf{Tokens} & \textbf{tok/char} \\
        \midrule
        GPT-2 & 50,257 & 355 & 0.4797 \\
        DeepSeek-V2 & 100,000 & 220 & 0.2973 \\
        DeepSeek-V3 & 128,000 & 182 & \textbf{0.2459} \\
        Qwen3-8B & 151,643 & -- & -- \\
        Ours & 52,877 & 185 & 0.2500 \\
        \bottomrule
    \end{tabular}
\end{table}

\textbf{Key observations:}
\begin{itemize}
    \item \texttt{quicksort} is tokenized as \texttt{quicks | ort} in GPT-2, but \texttt{quick | sort} in Ours -- more semantically meaningful
    \item DeepSeek-V3 achieves 182 tokens (best), Ours achieves 185 tokens (2nd best)
    \item Ours is within 2\% of DeepSeek-V3 despite having 2.5x fewer vocabulary entries
\end{itemize}

\textbf{Multi-language code comparison (tok/char):}

\begin{table}[h]
    \centering
    \begin{tabular}{lrrrr}
        \toprule
        \textbf{Language} & \textbf{GPT-2} & \textbf{DS-V3} & \textbf{Qwen3} & \textbf{Ours} \\
        \midrule
        Python & 0.4483 & 0.3241 & -- & 0.3379 \\
        JavaScript & 0.3421 & 0.2237 & -- & \textbf{0.2237} \\
        SQL & 0.4058 & \textbf{0.3043} & -- & 0.3261 \\
        \bottomrule
    \end{tabular}
\end{table}

Ours matches or approaches DeepSeek-V3's code tokenization quality across languages, confirming the effectiveness of our custom BPE training on code data.

\subsection{Multilingual Coverage: Beyond English and Chinese}

Qwen3-8B and DeepSeek-V3 invest their large vocabularies (151K and 128K) in broad multilingual coverage. To ensure a fair comparison, we evaluate tokenization across seven additional writing systems -- scripts where our tokenizer has no dedicated vocabulary budget.

\begin{table}[h]
    \centering
    \caption{Multilingual Tokenization: tok/char (lower is better)}
    \begin{tabular}{lrrrrr}
        \toprule
        \textbf{Script} & \textbf{GPT-2} & \textbf{DS-V3} & \textbf{Qwen3} & \textbf{Ours} & \textbf{Ours vs GPT-2} \\
        \midrule
        Russian (Cyrillic) & 1.0858 & \textbf{0.3729} & 0.4257 & 0.9373 & 1.16$\times$ better \\
        Arabic & 0.9116 & 0.3757 & \textbf{0.3702} & 0.7956 & 1.15$\times$ better \\
        Japanese & 1.3571 & 0.6073 & \textbf{0.5805} & 1.4198 & 0.96$\times$ \\
        Korean (Hangul) & 2.0982 & 0.6341 & \textbf{0.6252} & 2.2770 & 0.92$\times$ \\
        French & 0.3793 & \textbf{0.2530} & 0.2951 & 0.3679 & 1.03$\times$ better \\
        German & 0.3153 & \textbf{0.2267} & 0.2710 & 0.3104 & 1.02$\times$ better \\
        Unicode symbols & 2.0671 & 1.8537 & \textbf{1.1768} & 2.1280 & 0.97$\times$ \\
        \bottomrule
    \end{tabular}
\end{table}

\textbf{Key observations:}
\begin{itemize}
    \item \textbf{Qwen3 and DS-V3 dominate multilingual:} 7/7 wins across non-English/non-Chinese scripts, as expected given their larger vocabularies
    \item \textbf{vs GPT-2 (same scale):} Ours outperforms GPT-2 on Russian (+16\%), Arabic (+15\%), French (+3\%), German (+2\%); ties on Japanese and Unicode symbols; slightly worse on Korean (--8\%). This is consistent with Ours inheriting Qwen3-8B's multilingual merges through the hybrid merging pipeline.
    \item \textbf{Korean/Japanese gap:} Ours lacks dedicated Hangul/Kana tokens, resulting in character-level splitting (2.28 and 1.42 tok/char). This is expected given our 53K vocabulary budget.
    \item \textbf{Unicode symbols:} Qwen3's emoji/arrow/math symbol coverage is far superior (1.18 vs 2.13 for ours)
\end{itemize}

\textbf{Putting it in perspective:} Across all 13 test types (original 7 + 7 multilingual), the win count is:
\begin{itemize}
    \item Qwen3: 5 wins (Code, English, Unicode, Arabic, Japanese, Korean)
    \item DS-V3: 4 wins (Chinese Simp/Trad, Russian, French, German)
    \item GPT-2: 1 win (Numbers)
    \item Ours: 0 wins (but consistently beats GPT-2 across 10/13 tests)
\end{itemize}

\textbf{Putting it in perspective:} Across all 13 test types, the win count is:
\begin{itemize}
    \item Qwen3: 5 wins (English, Code, Unicode, Arabic, Japanese, Korean)
    \item DS-V3: 4 wins (Chinese Simp/Trad, Russian, French, German)
    \item GPT-2: 1 win (Numbers)
    \item Ours: 0 wins (but consistently beats GPT-2 across 10/13 tests)
\end{itemize}

\subsection{Tokenization Summary: Toward Usable SLM Agents}

From a \textbf{pragmatic, agent-oriented perspective}, tokenizer quality should be evaluated against the target use case, not against universal multilingual coverage. Our goal is to build a useful SLM agent operating with English, Chinese, and Code -- the three domains that matter for the agent paradigm we target.

\textbf{What matters for our agent:}
\begin{enumerate}
    \item \textbf{English:} Foundation for instruction following, tool use, reasoning
    \item \textbf{Chinese:} Requirement for bilingual agent deployment
    \item \textbf{Code:} Essential for tool calling, code generation, structured reasoning
\end{enumerate}

\begin{table}[h]
    \centering
    \caption{Agent-Relevant Tokenizer Performance Summary}
    \begin{tabular}{lrrrrr}
        \toprule
        \textbf{Domain} & \textbf{GPT-2} & \textbf{Ours} & \textbf{$\Delta$ (Ours vs GPT-2)} & \textbf{$\ln(V)$ Efficiency ($A$)} \\
        \midrule
        English & 0.2630 & 0.2400 & +9.6\% better & \textbf{2.610} (best) \\
        Chinese (Simp) & 2.0971 & 0.9045 & +56.9\% better & 9.841 \\
        Chinese (Trad) & 2.1442 & 1.0087 & +53.0\% better & 10.975 \\
        Code & 0.3817 & 0.3123 & +18.2\% better & 3.399 (2nd) \\
        \bottomrule
    \end{tabular}
\end{table}

\textbf{Conclusion:} Ours is not the best universal tokenizer -- Qwen3 and DS-V3 rightfully dominate on multilingual scripts where they invest 2-3$\times$ more vocabulary. However, for the \textbf{English + Chinese + Code} domain that defines our agent use case, Ours achieves:
\begin{itemize}
    \item \textbf{Best vocabulary efficiency} ($A=2.610$ on English, ranked \#1 across all tokenizers)
    \item \textbf{53-57\% better Chinese compression} than the same-scale GPT-2 baseline
    \item \textbf{18\% better Code compression} than GPT-2, within 2\% of DS-V3's 128K vocabulary
    \item \textbf{Solid multilingual transfer} from Qwen3-8B merges (beats GPT-2 on Russian, Arabic, French, German)
\end{itemize}

This represents a pragmatic trade-off: we accept limited coverage of non-target scripts (Japanese, Korean, Vietnamese) to achieve state-of-the-art tokenizer efficiency on the English-Chinese-Code triad that matters for our agent. The $\ln(V)$ scaling analysis confirms this is not just a large-vocabulary artifact but a genuine quality advantage in vocabulary utilization.

\section{Training Runs}
\label{sec:experiments}

\subsection{Overview}

\begin{table}[h]
    \centering
    \caption{Training Runs Summary}
    \begin{tabular}{lllrlll}
        \toprule
        \textbf{Run} & \textbf{Model} & \textbf{Scale} & \textbf{Steps} & \textbf{Loss} & \textbf{PPL} & \textbf{Status} \\
        \midrule
        tinystories & DenseSLM4 (dense) & 7M & 1,559 & 3.050 & 21.11 & Complete (smoke test) \\
        my\_model\_new\_tokenizer & DenseSLM4 (dense) & 105M & 122,837 & 2.143 & 8.52 & Complete \\
        denseslm4\_moe & DenseSLM4MoE v1 & $\sim$1.4B & 129,107 & 2.622 & 13.76 & Complete \\
        denseslm4\_moe\_continued & DenseSLM4MoE v1 + continued & $\sim$1.4B & +19,367 & 2.591 & 13.34 & Complete \\
        denseslm4\_moe\_v2 & DenseSLM4MoE v2 & \textbf{1.45B} & $\sim$55,900/190,967 & $\sim$2.4 & $\sim$11 & \textbf{In Progress (29\%)} \\
        moe\_low\_rank\_ablation & Small Trans+MoE & 25M & 1,000 & 2.033 & 7.63 & Complete (ablation) \\
        \bottomrule
    \end{tabular}
\end{table}

\subsection{DenseSLM4MoE v2 Training Progress}

\begin{itemize}
    \item \textbf{Start date:} $\sim$May 20, 2026
    \item \textbf{Current step:} $\sim$55,900 of 190,967 planned (29\%)
    \item \textbf{Checkpoints saved:} ckpt-30000, ckpt-40000, ckpt-50000
    \item \textbf{Loss trajectory:} Step 0: $\sim$3.5 $\rightarrow$ Step 50K: $\sim$2.4 (rapid improvement in first half)
    \item \textbf{Remaining:} $\sim$135K steps estimated $\sim$50 additional hours at current speed
\end{itemize}

\section{Evaluation}
\label{sec:evaluation}

\subsection{Current Results}

\subsubsection{DenseSLM4MoE v1 (1.4B params, 22K vocab)}

\begin{table}[h]
    \centering
    \caption{DenseSLM4MoE v1 LightEval Results}
    \begin{tabular}{llrr}
        \toprule
        \textbf{Task} & \textbf{Metric} & \textbf{0-shot} & \textbf{5-shot} \\
        \midrule
        C-Eval (avg, 43 subjects) & Accuracy & 0.217 & \textbf{0.281} \\
        TruthfulQA \citep{truthfulqa} & MC1 & 0.203 & --- \\
        TruthfulQA & MC2 & 0.440 & --- \\
        \bottomrule
    \end{tabular}
\end{table}

\subsubsection{DenseSLM4 (Dense, 105M params)}

\begin{table}[h]
    \centering
    \caption{DenseSLM4 LightEval Results}
    \begin{tabular}{llr}
        \toprule
        \textbf{Task} & \textbf{Metric} & \textbf{Score} \\
        \midrule
        GSM8K \citep{gsm8k} (5-shot) & Extractive Match & 0.076\% \\
        TruthfulQA (0-shot) & MC1 / MC2 & 0.240 / 0.430 \\
        \bottomrule
    \end{tabular}
\end{table}

\subsection{Analysis: Gap to Useful Agent Level}

Current results on standard benchmarks are \textbf{modest}, which is expected for models at this scale without extensive RLHF:

\begin{table}[h]
    \centering
    \caption{Capability Gap Analysis}
    \begin{tabular}{lllr}
        \toprule
        \textbf{Capability} & \textbf{Target} & \textbf{Current} & \textbf{Gap} \\
        \midrule
        Chinese reasoning (C-Eval 5-shot) \citep{ceval} & $\ge$35\% & 28.1\% & Needs more Chinese data + RLHF \\
        Truthfulness (TruthfulQA MC1) & $\ge$40\% & 20-24\% & Needs RLHF \\
        Math reasoning (GSM8K) & $\ge$20\% & $\sim$0\% & Needs more math data + CoT \\
        Code generation (HumanEval) \citep{humaneval} & $\ge$20\% & Not evaluated & Planned for v2 \\
        \bottomrule
    \end{tabular}
\end{table}

\subsection{What DenseSLM4MoE v2 Has That Helps}

\begin{enumerate}
    \item \textbf{Larger vocab (53K vs 22K):} Better tokenization of Chinese and code
    \item \textbf{More training steps:} v1 used 129K steps; v2 will reach 191K steps
    \item \textbf{Better tokenizer compression:} $\sim$1.5 tokens/char (v2) vs $\sim$1.3 (v1)
    \item \textbf{WCL schedule:} Better LR management for final capability push
\end{enumerate}

\subsection{Per-Subject C-Eval Analysis (v1, 5-shot)}

\textbf{Strongest subjects ($>$40\%):} College Physics (52.6\%), High School Politics (52.6\%), Computer Network (52.6\%)

\textbf{Weakest subjects ($<$20\%):} Middle School History (13.6\%), Physician (16.3\%), Clinical Medicine (18.2\%)

This suggests the model has better physics/CS knowledge than social sciences --- consistent with English-heavy pretraining dominating the Chinese educational data.

\subsection{Architecture Comparison with GPT-2}

\begin{table}[h]
    \centering
    \caption{DenseSLM4MoE v2 vs GPT-2 XL}
    \begin{tabular}{lll}
        \toprule
        \textbf{Property} & \textbf{GPT-2 XL (1.5B)} & \textbf{DenseSLM4MoE v2 (1.45B)} \\
        \midrule
        Architecture & Transformer (dense) & Hybrid MLA + Mamba2 + MoE \\
        Attention & Full MHA (16 heads) & MLA compressed (4 KV heads, 512-dim) \\
        FFN & Dense (intermediate=4$\times$) & Sparse MoE (64 experts, top-2) \\
        KV Cache & Full (O(N) growth) & Compressed (constant 512-dim/head) \\
        Layer count & 48 & 48 \\
        Hidden size & 1600 & 768 \\
        Vocab size & 50,257 & 52,877 \\
        Active params/pass & $\sim$1.5B & $\sim$768M (\textbf{2$\times$ more efficient}) \\
        Total params & 1.5B & 1.45B \\
        \bottomrule
    \end{tabular}
\end{table}

\section{Conclusion}
\label{sec:conclusion}

DenseSLM4MoE demonstrates a practical hybrid architecture achieving GPT-2-scale parameters (1.45B) with significantly improved active-parameter efficiency ($\sim$768M active vs 1.45B total). The combination of MLA (compressed KV), Mamba2 (efficient SSM layers), and aux-loss-free MoE provides a strong foundation for agent workloads at small scale.

\subsection{Current Status}
\begin{itemize}
    \item DenseSLM4MoE v2 pretraining is complete (190{,}967 steps, one epoch over $\sim$6.3B tokens), reaching an evaluation perplexity of 11.49
    \item Architecture is validated through ablation and v1 experiments
    \item Training is stable with no loss spikes or divergence
\end{itemize}

\subsection{What's Needed to Reach ``Useful Agent Level''}

Concrete next steps are detailed in Section~\ref{sec:futurework}.

\subsection{Limitations}
\begin{itemize}
    \item Current benchmarks (C-Eval 28.1\%, TruthfulQA 20\%) are below the threshold for truly useful agents
    \item No RLHF or DPO has been applied yet --- pure next-token prediction pretraining
    \item GSM8K near-random suggests mathematical reasoning needs targeted improvement
    \item Chinese performance bottlenecked by data quality, not architecture
\end{itemize}

\subsection*{Acknowledgments}

We thank the \texttt{fla} \citep{fla} and \texttt{mamba\_ssm} teams for their excellent open-source implementations of Mamba2 and MLA. We also thank the HuggingFace team for LightEval \citep{lighteval}.

\subsection*{Reproducibility}

All model configurations, training scripts, and evaluation pipelines are available at:
\begin{verbatim}
https://github.com/ZiceWang/denseslm4
\end{verbatim}

Training can be launched via:
\begin{verbatim}
python -m denseslm4.train_moe --output_dir runs/denseslm4_moe_v2 ...
\end{verbatim}

Evaluation:
\begin{verbatim}
python -m lighteval --model_args "pretrained=./runs/denseslm4_moe/final_model,..." ...
\end{verbatim}

\section{Future Work}
\label{sec:futurework}

The present results characterise a first working point of the hybrid design
(1.45B total / $\sim$768M active parameters, $\sim$6.3B training tokens). The
following steps are needed before any claim about ``useful agent'' capability, or
about the architectural choices themselves, can be substantiated.

\paragraph{Scaling pretraining compute.}
Training consumed roughly $6.3$B tokens (one epoch over $6.1$M $\times$ 1024-token
chunks), which is well below compute-optimal for a model of this size. Scaling to
$10^2$--$10^3$B tokens is expected to be the single largest source of improvement and
is the natural next run.

\paragraph{Matched-baseline comparison.}
The current evaluation does not isolate the contribution of the hybrid architecture.
A controlled study at an identical token budget should compare against (i) a dense
Transformer of matched active parameters and (ii) a public same-scale base model
(e.g.\ Pythia-1.4B) at matched tokens, before attributing any gain to MLA, Mamba2, or
the aux-loss-free MoE.

\paragraph{Post-training.}
So far only next-token-prediction pretraining and a preliminary SFT stage have been
run. Supervised instruction tuning followed by preference optimisation (DPO/RLHF)
should be evaluated end-to-end, since the current gap is dominated by missing
preference learning rather than by the architecture.

\paragraph{Evaluation coverage.}
Beyond C-Eval, TruthfulQA, and GSM8K, the agent setting requires code generation
(HumanEval/MBPP), long-context retrieval, and multi-turn tool-use tasks. Every reported
score should be accompanied by the token budget and by matched-baseline deltas.

\paragraph{Architecture ablations.}
The routed-expert count and top-$k$, the MLA KV-compression rank, and the
Mamba2/attention interleaving ratio should be swept under a fixed token budget, so that
each component's contribution is measured rather than assumed.

\paragraph{Long context.}
The current context window is 1024 tokens. Extending it via position interpolation and
continued pretraining is a prerequisite for the document- and tool-heavy workloads the
architecture targets.

\paragraph{Inference efficiency.}
The constant-size KV cache and SSM state, together with the TileLang kernel path, should
be benchmarked on throughput and peak memory, in order to quantify the active-parameter
efficiency claim on realistic agent workloads.


\bibliographystyle{plainnat}
\bibliography{references}

\end{document}
